\documentclass[12pt]{article}
\usepackage[margin=1in]{geometry}
\usepackage{amsmath, amssymb}
\usepackage{enumitem, array}

\usepackage{boxedminipage}
\usepackage{fancyhdr, actuarialangle}
\usepackage{ragged2e, comment}


\newenvironment{problem}[1][]{%
  \bigskip% put space before problem statement box %
  \noindent\begin{boxedminipage}{\columnwidth}%
  \def\@tempa{#1}%
  \ifx\@tempa\empty\else%
    \textbf{#1}\hspace{0.5em}%
  \fi%
}{%
  \end{boxedminipage}%
}


\begin{document}

\begin{center}
    \Large \textbf{MATH 2790 Test 2 Fall 2025} \\[1em]
\end{center}

\vspace{.25cm}

\noindent Name: \underline{\hspace{8cm}} \hspace{.25cm} Auburn ID: \underline{\hspace{4cm}}



\section*{Instructions}
\begin{enumerate}[label=\arabic*.]
\item Print your full name and your banner ID on your exam. Then finish reading these instructions, and sign the bottom of this page.
\item All electronic equipment (cell phones, Apple Watch, AirPod, etc) are \textbf{NOT ALLOWED} during the exam time.
\item You are allowed to use up to 2 calculators from the following list: BA-35, BA II Plus, BA II Plus Professional, TI-30Xa, TI-30X II (IIS solar or IIB battery), TI-30XS MultiView (or XB battery).
\item No talking is allowed during the exam.
\item Unless otherwise indicated, \textbf{SHOW ALL YOUR WORK}. If no work is shown, no partial credit can be awarded. Your work and answer need to be accurate and relevant to receive points.
\item You will be given exactly 50 minutes for this exam.
\item Any student not following the above instructions nor behaving according to the above instructions during the exam may have their exam confiscated and points deducted.
\end{enumerate}

I have read and fully understand all of the above instructions. \\[2em]
Signature: \hrulefill


\vspace{.75cm}
Do not write in the table below

% scoring table
\begin{center}
\begin{tabular}{|*{3}{>{\centering\arraybackslash}p{.1\linewidth}|}}
    \hline
    \textbf{Problem} & \textbf{Points} & \textbf{Score} \\
    \hline
    1 & 15 &  \\
    \hline
    2 & 15 &  \\
    \hline
    3 & 15 &  \\
    \hline
    4 & 15 &  \\
    \hline
    5 & 16 &  \\
    \hline
    6 & 12 &  \\
    \hline
    7 & 12 &  \\
    \hline \hline
    \textbf{Total} & 100 &  \\
    \hline
\end{tabular}
\end{center}

\newpage

%%%%% PROBLEM 1 %%%%%
\newpage
\begin{problem}[1. (15 points)]
Jared buys an increasing annuity. Jared receives \$3 at the end of the first month, \$5 at the end of the second month, and for each month thereafter the payment increases by \$2. The amount of the last payment is \$41. The nominal interest rate is 6\% convertible monthly. Calculate the value of the annuity at the time of his first payment.
\end{problem}

\newpage
%%%%% PROBLEM 2 %%%%%
\newpage
\begin{problem}[2. (15 points)]
Cindy, on her 40th birthday, received a perpetuity that will pay \$20{,}000 annually starting immediately. The annual interest rate is 6\% up until her 65th birthday and 3\% thereafter. Find the value of the perpetuity on the day Cindy received it. Round your answer to two decimal places.
\end{problem}

\newpage

%%%%% PROBLEM 3 %%%%%
\newpage
\begin{problem}[3. (15 points)]
As part of a new initiative, at the end of each year, a city government will allocate funds every year for upkeep on a public transit system for the next 20 years. Prior to this project, annual costs were \$3{,}000 for year 2023 and are expected to increase by 6\% per year due to inflation. The first payment of the 20 year project will be made on December 31, 2024.  

\bigskip

Calculate the value of the obligation using an annual effective interest rate of 4\% on December 31, 2024. Round your answer to four decimal places.
\end{problem}

\newpage

%%%%% PROBLEM 4 %%%%%
\newpage
\begin{problem}[4. (15 points)]
An annuity is continuously varying and payable for 10 years. The rate of payment at time \(t\) is \(t\), and the force of interest is \(\delta_t = 0.05\). Find the present value of this annuity. Round your answer to two decimal places.
\end{problem}

\newpage

%%%%% PROBLEM 5 %%%%%
\newpage
\begin{problem}[5. (16 points)]
Bob purchases a really really nice boat. The price of the boat is 200,000 and he makes a down payment of 80,000. He finances the balance with monthly payments for $5$ years with an annual interest rate of $8.94\%$.
\begin{enumerate}
    \item[(a)] If Bob makes level payments at the end of each month, what is the value of the payment? Round your answer to $4$ decimal places. (6 points)
    \item[(b)] Bob does not like loose change (pennies, nickels, etc.) so he rounds up and makes the payments from (a) rounded up to the nearest dollar (except for the final payment). Find the outstanding loan balance right after his 45th payment. Round to $2$ decimal places. (10 points)
\end{enumerate}
\end{problem}

\newpage

%%%%% PROBLEM 6 %%%%%
\newpage
\begin{problem}[6. (12 points)]
You plan to support \$75{,}000 for your child’s college tuition in the future.  
Thus, you decide to make quarterly deposits for 10 years in a savings account.  
The final deposit is made 3 months prior to the college payment of \$75{,}000.  
Assume the savings account earns the annual interest rate \(i\).  

\bigskip

Let \(X\) be the amount of each quarterly deposit. Set up the equation of \(X\) using actuarial notation and annual effective interest rate \(i\).  
\textbf{No need to solve for \(X\)}.
\end{problem}

\newpage

%%%%% PROBLEM 7 %%%%%
\newpage
\begin{problem}[7. (12 points)]
Payments of 1 are made at the end of year 4, 8, 12, 16, and 20 with annual effective interest rate of \( i \).

\bigskip

Circle your answer(s) that express the present value of the annuity and \textbf{show work to justify your answer(s)}.

\bigskip

\begin{tabbing}
\hspace{0.6cm}\=\hspace{1cm}\=\kill
(a) \> {$\ddot{a}_{\actuarialangle{5} i} v^4$} \\[1.5ex]
(b) \> {$\dfrac{a_{\actuarialangle{20} i}}{\ddot{a}_{\actuarialangle{4} i}}$} \\[2ex]
(c) \> {$\dfrac{a_{\actuarialangle{20} i}}{s_{\actuarialangle{4} i}}$} \\[2ex]
(d) \> {$\dfrac{a_{\actuarialangle{20} i}}{\ddot{s}_{\actuarialangle{4} i}}$}
\end{tabbing}
\end{problem}

\newpage

%%%% PROBLEM 8 %%%%%
\newpage
\begin{problem}[8. (+1 Bonus point)]
Give me two numbers the first between $0$ and $\pi \approx 3.1415$ and the second between $0$ and $2\pi \approx 6.28318$ (You can be creative and the numbers can have as many decimals as you wish!)


\vspace{1cm}

Example: $(.375375375, 1.111111)$
\end{problem}


\end{document}
